\chapter{Quantization, the Spectrum and the Critical Dimension}\label{ch:quantum}

\Cref{ch:classical} described a classical relativistic string: $D$ free fields $X^\mu(\sigma,\tau)$
on the worldsheet, subject to constraints. This chapter asks what that system becomes in quantum
mechanics. Which particles does one quantized string describe? What are their masses and spins?
And why does the answer depend on the number $D$ of spacetime dimensions, so that a consistent
bosonic string exists only for $D=26$? A student should care for two reasons. First, the
spectrum derived here contains a massless spin-two particle, and this is the reason string
theory is a theory of gravity. It also contains the two other massless fields, $B_{\mu\nu}$ and
the dilaton $\Phi$, from which the whole of T-duality and double field theory in Parts II and
III is built. Second, the mass formula obtained here, $M^2=\frac{4}{\ap}(N-1)$, is the
$R\to\infty$ limit of the formula with momentum and winding in \cref{ch:circle}; every later
statement about dual scales is a statement about how this formula changes on a circle.
The chapter closes by stating exactly which part of the ``$D=26$'' story has been checked in
Lean~4: the arithmetic of central charges, and nothing more.

Conventions are those of \cref{ch:classical}: $\hbar=c=1$, mostly-plus Minkowski metric
$\eta_{\mu\nu}$, $\ap$ of dimension length$^2$, tension $T=1/(2\pi\ap)$, closed strings with
$\sigma\sim\sigma+2\pi$. A dot between two vectors is the Minkowski product,
$A\cdot B=\eta_{\mu\nu}A^\mu B^\nu$.

%=====================================================================================
\section{Canonical quantization: from fields to oscillators}\label{sec:quantum:canonical}

\subsection{Two roads}

The string action has a gauge symmetry (reparametrizations and Weyl rescalings of the
worldsheet), and a gauge theory can be quantized in two ways \source{Tong §2, ll.~1683--1707}.
\begin{itemize}[leftmargin=*]
\item \emph{Covariant quantization}\index{covariant quantization}: quantize all $D$ fields
$X^\mu$, including the time $X^0$, and afterwards impose the constraints as conditions on
states. This is the analogue of the Gupta--Bleuler treatment of electrodynamics in Lorentz
gauge. Lorentz invariance is manifest; positivity of the Hilbert space is not.
\item \emph{Lightcone quantization}\index{lightcone quantization}: solve the constraints
classically, find the independent degrees of freedom, and quantize only those. This is the
analogue of Coulomb gauge. Positivity is manifest; Lorentz invariance is not.
\end{itemize}
If both are carried out correctly they agree. We follow Tong: a short look at the first
road, to see where the difficulties are, and then the second road to the end.

\subsection{Commutation relations for the modes}

In conformal gauge the action is that of $D$ free scalar fields,
$S=-\frac{1}{4\pi\ap}\int\dd^2\sigma\,\partial_\alpha X\cdot\partial^\alpha X$, and the
constraints are
\begin{equation}\label{eq:quantum:constraints}
  \dot X\cdot X'=0,\qquad \dot X^2+X'^2=0 ,
\end{equation}
where a dot is $\partial_\tau$ and a prime is $\partial_\sigma$
\source{Tong §2.1, ll.~1708--1713}. With worldsheet lightcone coordinates
$\sigma^\pm=\tau\pm\sigma$ the general periodic solution of the wave equation
$\partial_+\partial_-X^\mu=0$ is $X^\mu=X^\mu_L(\sigma^+)+X^\mu_R(\sigma^-)$ with
\begin{equation}\label{eq:quantum:modes}
\begin{aligned}
 X_L^\mu(\sigma^+)&=\tfrac12x^\mu+\tfrac12\ap p^\mu\sigma^+
   +\ii\sqrt{\tfrac{\ap}{2}}\sum_{n\neq0}\frac1n\,\tilde\alpha^\mu_n\,\ee^{-\ii n\sigma^+},\\
 X_R^\mu(\sigma^-)&=\tfrac12x^\mu+\tfrac12\ap p^\mu\sigma^-
   +\ii\sqrt{\tfrac{\ap}{2}}\sum_{n\neq0}\frac1n\,\alpha^\mu_n\,\ee^{-\ii n\sigma^-},
\end{aligned}
\end{equation}
\source{Tong §1.4, ll.~1502--1563}. Here $x^\mu$ and $p^\mu$ are the position and momentum of
the centre of mass, $\alpha_n^\mu$ are the amplitudes of right-moving waves,
$\tilde\alpha_n^\mu$ those of left-moving waves, and reality of $X^\mu$ requires
$\alpha_{-n}^\mu=(\alpha_n^\mu)^*$. One also defines
$\alpha_0^\mu=\tilde\alpha_0^\mu=\sqrt{\ap/2}\,p^\mu$.

The momentum conjugate to $X^\mu$ is $\Pi_\mu=\partial\mathcal L/\partial\dot X^\mu
=\dot X_\mu/(2\pi\ap)$. Quantization replaces fields by operators with the equal-time
commutators\index{canonical quantization}
\begin{equation}\label{eq:quantum:ccr}
 [X^\mu(\sigma,\tau),\Pi_\nu(\sigma',\tau)]=\ii\,\delta(\sigma-\sigma')\,\delta^\mu{}_\nu,
 \qquad [X^\mu,X^\nu]=[\Pi_\mu,\Pi_\nu]=0 ,
\end{equation}
and these are equivalent to
\begin{equation}\label{eq:quantum:osc}
 [x^\mu,p_\nu]=\ii\,\delta^\mu{}_\nu,\qquad
 [\alpha^\mu_n,\alpha^\nu_m]=[\tilde\alpha^\mu_n,\tilde\alpha^\nu_m]
   =n\,\eta^{\mu\nu}\,\delta_{n+m,0},
\end{equation}
with all other commutators zero \source{Tong §2.1, ll.~1714--1742}.

It is worth checking once that \cref{eq:quantum:osc} reproduces \cref{eq:quantum:ccr},
because the check fixes every normalization in \cref{eq:quantum:modes}. At $\tau=0$ one has
$\ee^{-\ii n\sigma^-}=\ee^{\ii n\sigma}$ and $\ee^{-\ii n\sigma^+}=\ee^{-\ii n\sigma}$, and
\[
 \dot X^\nu(\sigma')=\ap p^\nu+\sqrt{\tfrac{\ap}{2}}\sum_{m\ne0}
   \big(\alpha^\nu_m\ee^{\ii m\sigma'}+\tilde\alpha^\nu_m\ee^{-\ii m\sigma'}\big).
\]
The zero modes give $[x^\mu,\ap p^\nu]=\ii\ap\eta^{\mu\nu}$. In the oscillator part the
commutator $[\alpha^\mu_n,\alpha^\nu_m]=n\eta^{\mu\nu}\delta_{n+m,0}$ cancels the factor $1/n$
and sets $m=-n$:
\[
 \ii\,\frac{\ap}{2}\sum_{n\ne0}\frac1n\,n\,\eta^{\mu\nu}
   \big(\ee^{\ii n(\sigma-\sigma')}+\ee^{-\ii n(\sigma-\sigma')}\big)
 =\ii\ap\,\eta^{\mu\nu}\sum_{n\neq0}\ee^{\ii n(\sigma-\sigma')} .
\]
Add the two pieces and use the Fourier series of the periodic delta function,
$\sum_{n\in\Z}\ee^{\ii n(\sigma-\sigma')}=2\pi\delta(\sigma-\sigma')$. The result is
\[
 [X^\mu(\sigma),\dot X^\nu(\sigma')]=2\pi\ii\ap\,\eta^{\mu\nu}\,\delta(\sigma-\sigma') ,
\]
and dividing by $2\pi\ap$ gives \cref{eq:quantum:ccr}. The factor $n$ on the right of
\cref{eq:quantum:osc} is therefore forced by the $1/n$ in the mode expansion.

\subsection{Harmonic oscillators and the Fock space}

For a single direction $\mu$ and $n>0$ define
\begin{equation}\label{eq:quantum:ladder}
  a_n=\frac{\alpha_n}{\sqrt n},\qquad a_n^\dagger=\frac{\alpha_{-n}}{\sqrt n}
  \qquad\Longrightarrow\qquad [a_n,a_m^\dagger]=\delta_{nm}
\end{equation}
\source{Tong §2.1, ll.~1743--1757}. Each spatial field $X^i$ is therefore two infinite
towers of harmonic oscillators: $\alpha_n$ with $n>0$ annihilates a right-moving quantum of
worldsheet frequency $n$ and $\alpha_{-n}$ creates one; the $\tilde\alpha_n$ do the same for
left-movers.

\begin{physicsbox}[Physical picture: a violin string with infinitely many notes]
A classical closed string vibrates in normal modes labelled by an integer $n$ (the number of
wavelengths around the string), a direction $\mu$ and a sense of travel (left or right).
Quantization turns each normal mode into a harmonic oscillator whose excitations come in
integer quanta. Seen from far away, a small vibrating loop looks like a point particle, and
the vibrational energy appears as the rest mass of that particle. One string, excited in
different ways, is therefore an infinite list of particle species
\source{Tong §2.1, ll.~1790--1800}.
\end{physicsbox}

The vacuum of the oscillators is defined by
\begin{equation}\label{eq:quantum:vacuum}
 \alpha^\mu_n\ket{0;p}=\tilde\alpha^\mu_n\ket{0;p}=0\quad(n>0),\qquad
 \hat p^\mu\ket{0;p}=p^\mu\ket{0;p},
\end{equation}
where the label $p$ is needed because the zero modes $x^\mu,p^\mu$ act as in the quantum
mechanics of a free particle. This state is the ground state of \emph{one string} with
centre-of-mass momentum $p$. It is not the vacuum of spacetime, which contains no string at
all \source{Tong §2.1, ll.~1758--1782}. The Fock space\index{Fock space} is spanned by the
states obtained by acting on $\ket{0;p}$ with any finite number of creation operators
$\alpha^\mu_{-n},\tilde\alpha^\mu_{-n}$, $n>0$.

%=====================================================================================
\section{Negative norms and the Virasoro constraints}\label{sec:quantum:virasoro}

\subsection{The problem with the time direction}

The Fock space just built is not a Hilbert space. Because $\eta^{00}=-1$, the commutator
in \cref{eq:quantum:osc} for the time direction reads $[\alpha^0_n,\alpha^0_{-n}]=-n$, so
\begin{equation}\label{eq:quantum:negnorm}
 \big\|\alpha^0_{-1}\ket{0;p}\big\|^2=\bra{0;p'}\alpha^0_1\alpha^0_{-1}\ket{0;p}
   =\bra{0;p'}[\alpha^0_1,\alpha^0_{-1}]\ket{0;p}=-\,\delta^D(p-p') .
\end{equation}
States with an odd number of timelike oscillators have negative norm; such states are called
\emph{ghosts}\index{ghost (negative-norm state)} \source{Tong §2.1.1, ll.~1802--1822}. A
negative norm would mean a negative probability, so these states must not appear in any
physical process. In electrodynamics the same problem occurs for the timelike photon, and
the gauge-fixing condition removes it. In string theory the constraints of
\cref{eq:quantum:constraints} must do this job.

\begin{pitfallbox}[Common pitfall: two meanings of ``ghost'']
In this chapter a ghost is a \emph{state of negative norm}. In \cref{ch:ghosts} the word
also denotes the anticommuting Faddeev--Popov fields $b,c$ introduced when the path integral
is gauge-fixed. The two are related, since the Faddeev--Popov fields are what cancels the
unphysical polarizations, but they are different objects. The Lean file discussed in
\cref{sec:quantum:lean} is about the central charge of the second kind.
\end{pitfallbox}

\subsection{The constraints in modes}

In lightcone worldsheet coordinates the constraints of \cref{eq:quantum:constraints} read
$(\partial_+X)^2=(\partial_-X)^2=0$. From \cref{eq:quantum:modes},
$\partial_-X^\mu=\sqrt{\ap/2}\,\sum_{n\in\Z}\alpha^\mu_n\ee^{-\ii n\sigma^-}$, where the $n=0$
term is the definition of $\alpha_0^\mu$. Squaring,
\begin{equation}\label{eq:quantum:Ln}
 (\partial_-X)^2=\ap\sum_{n\in\Z}L_n\,\ee^{-\ii n\sigma^-},\qquad
 L_n=\frac12\sum_{m\in\Z}\alpha_{n-m}\cdot\alpha_m ,
\end{equation}
and similarly $\tilde L_n$ from the left-movers \source{Tong §1.4.1, ll.~1565--1630}.
Classically the constraints are $L_n=\tilde L_n=0$ for all $n$. The $L_n$ are the
\emph{Virasoro generators}\index{Virasoro generators}; their algebra, and the central charge
that appears in it, are the subject of \cref{ch:cft}.

The $n=0$ constraint is special because it contains
$\frac12\alpha_0^2=\frac{\ap}{4}p^2=-\frac{\ap}{4}M^2$, where $M$ is the rest mass of the
string seen as a particle, $p_\mu p^\mu=-M^2$. Classically
$L_0=\frac{\ap}{4}p^2+\sum_{n>0}\alpha_{-n}\cdot\alpha_n=0$ gives
\begin{equation}\label{eq:quantum:classmass}
 M^2=\frac4\ap\sum_{n>0}\alpha_{-n}\cdot\alpha_n
    =\frac4\ap\sum_{n>0}\tilde\alpha_{-n}\cdot\tilde\alpha_n
\end{equation}
\source{Tong §1.4.1, ll.~1650--1680}. The mass is the vibrational energy. The second equality,
from $\tilde L_0=0$, holds because $\alpha_0=\tilde\alpha_0$: left- and right-movers share the
same centre-of-mass momentum. It is called \emph{level matching}\index{level matching}.

\subsection{Imposing the constraints on states}

In the quantum theory one cannot demand $L_n\ket{\rm phys}=0$ for all $n$; as in
Gupta--Bleuler quantization one demands only that matrix elements between physical states
vanish. Since $L_n^\dagger=L_{-n}$, it suffices that
\begin{equation}\label{eq:quantum:phys}
 L_n\ket{\rm phys}=\tilde L_n\ket{\rm phys}=0\qquad (n>0)
\end{equation}
\source{Tong §2.1.2, ll.~1823--1846}. For $n\neq0$ the two factors in each term of
\cref{eq:quantum:Ln} commute, so $L_n$ is unambiguous. For $n=0$ they do not: $\alpha_{-m}$
and $\alpha_m$ fail to commute, and different orderings differ by a constant. One fixes an
ordering, the normal ordering with annihilation operators on the right,
\begin{equation}\label{eq:quantum:L0}
 L_0=\frac12\alpha_0^2+\sum_{m=1}^\infty\alpha_{-m}\cdot\alpha_m ,\qquad
 \tilde L_0=\frac12\tilde\alpha_0^2+\sum_{m=1}^\infty\tilde\alpha_{-m}\cdot\tilde\alpha_m ,
\end{equation}
and parametrizes the ignorance by a constant $a$:
\begin{equation}\label{eq:quantum:L0a}
 (L_0-a)\ket{\rm phys}=(\tilde L_0-a)\ket{\rm phys}=0 .
\end{equation}
Inserting $\frac12\alpha_0^2=-\frac\ap4M^2$ into \cref{eq:quantum:L0}, these two conditions
read
\begin{equation}\label{eq:quantum:massL0a}
 M^2=\frac4\ap\Big(\sum_{m\ge1}\alpha_{-m}\cdot\alpha_m-a\Big)
    =\frac4\ap\Big(\sum_{m\ge1}\tilde\alpha_{-m}\cdot\tilde\alpha_m-a\Big)
\end{equation}
\source{Tong §2.1.2, ll.~1847--1912}. The constant $a$\index{normal-ordering constant} is not
a harmless convention: it shifts the mass of every state. Determining $a$, together with $D$,
is the main task of this chapter.

\begin{example}[The first excited level, covariantly]\label{ex:quantum:level1}
This example shows why $a=1$ is special before we derive it. For brevity keep only the
right-movers (the left-movers behave identically) and consider
$\ket{\zeta}=\zeta_\mu\alpha^\mu_{-1}\ket{0;p}$ with a polarization vector $\zeta$. The
calculation is a short exercise with \cref{eq:quantum:osc}; the argument is standard but is
not spelled out in this book's source library, where Tong refers to Green, Schwarz and
Witten for it \source{Tong §2.1.2, ll.~1920--1931}.

\emph{Mass shell.} $\sum_{m\ge1}\alpha_{-m}\cdot\alpha_m\ket\zeta=\ket\zeta$, so
\cref{eq:quantum:massL0a} gives $M^2=-p^2=4(1-a)/\ap$.

\emph{Constraint.} In $L_1=\alpha_0\cdot\alpha_1+\alpha_{-1}\cdot\alpha_2+\dots$ only the first
term acts non-trivially:
$L_1\ket\zeta=\alpha_0^\mu\zeta^\nu[\alpha_{1\mu},\alpha_{-1\,\nu}]\ket{0;p}
=\sqrt{\ap/2}\,(p\cdot\zeta)\ket{0;p}$. So \cref{eq:quantum:phys} requires $p\cdot\zeta=0$: the
polarization is transverse to the momentum. $L_n$ with $n\ge2$ annihilates $\ket\zeta$
automatically.

\emph{Norm.} $\langle\zeta|\zeta\rangle\propto\zeta^*\cdot\zeta$. Three cases:
\begin{itemize}[leftmargin=*]
\item $a<1$: $M^2>0$. In the rest frame $p=(M,0,\dots,0)$, so $p\cdot\zeta=0$ removes
$\zeta^0$. The $D-1$ remaining states have positive norm: a massive vector.
\item $a>1$: $M^2<0$, $p$ is spacelike, say $p=(0,\dots,0,|M|)$. Then $\zeta=(1,0,\dots,0)$
is allowed and has $\zeta\cdot\zeta=-1$: a physical ghost. The theory is inconsistent.
\item $a=1$: $M^2=0$. Take $p=(k,0,\dots,0,k)$. The condition $p\cdot\zeta=0$ allows the
$D-2$ transverse $\zeta$, of positive norm, and also $\zeta\propto p$, of norm
$p\cdot p=0$. This null state is orthogonal to all physical states and decouples. It
expresses the gauge invariance $\zeta_\mu\to\zeta_\mu+\lambda p_\mu$ of a massless vector.
\end{itemize}
The boundary value $a=1$ is where the theory is as rich as possible without ghosts: a massless
particle with $D-2$ polarizations and a gauge symmetry. For the closed string, with both
$\alpha_{-1}$ and $\tilde\alpha_{-1}$ acting, the same computation gives a polarization tensor
$\zeta_{\mu\nu}$ with $p^\mu\zeta_{\mu\nu}=\zeta_{\mu\nu}p^\nu=0$ and the invariance
$\zeta_{\mu\nu}\to\zeta_{\mu\nu}+p_\mu\xi_\nu+\xi'_\mu p_\nu$, which in its symmetric part is
the linearized diffeomorphism $h_{\mu\nu}\to h_{\mu\nu}+\partial_\mu\xi_\nu+\partial_\nu\xi_\mu$
\source{Tong §2.3.2, ll.~2725--2770}.
\end{example}

The full covariant analysis shows that the absence of ghosts at all levels requires $a=1$
and $D=26$ \TL\ \source{Tong §2.1.2, ll.~1920--1931}. We shall obtain the same values on the
second road.

%=====================================================================================
\section{Lightcone gauge: solving the constraints}\label{sec:quantum:lightcone}

\subsection{The residual gauge symmetry}

Conformal gauge $g_{\alpha\beta}=\eta_{\alpha\beta}$ does not use up all the gauge freedom.
The flat worldsheet metric is $\dd s^2=-\dd\sigma^+\dd\sigma^-$, so a reparametrization
\begin{equation}\label{eq:quantum:residual}
 \sigma^+\to\tilde\sigma^+(\sigma^+),\qquad \sigma^-\to\tilde\sigma^-(\sigma^-)
\end{equation}
only multiplies the metric by a function, which a Weyl rescaling removes
\source{Tong §2.2, ll.~1932--1962}. This has a consequence for the counting of degrees of
freedom. The solutions $X_L^\mu(\sigma^+)+X_R^\mu(\sigma^-)$ are $2D$ functions of one
variable. The constraints $(\partial_\pm X)^2=0$ remove two, and the freedom of
\cref{eq:quantum:residual} shows that two more are pure gauge. What remains is
\begin{equation}\label{eq:quantum:count}
 2D-2-2=2(D-2)\ \text{functions},
\end{equation}
the left- and right-moving \emph{transverse} oscillations of the string
\source{Tong §2.2, ll.~1963--1985}. A relativistic string has no longitudinal vibrations: a
motion of the string along itself is a relabelling of $\sigma$ and not a physical motion.

\subsection{The gauge choice}

Introduce spacetime lightcone coordinates\index{lightcone coordinates}
\begin{equation}\label{eq:quantum:lc}
 X^\pm=\frac{1}{\sqrt2}\big(X^0\pm X^{D-1}\big),\qquad
 \dd s^2=-2\,\dd X^+\dd X^-+\sum_{i=1}^{D-2}\dd X^i\dd X^i,
\end{equation}
so that $A\cdot B=-A^+B^--A^-B^++A^iB^i$ and $A_\pm=-A^\mp$
\source{Tong §2.2.1, ll.~1992--2035}. The choice singles out one spatial direction, and
Lorentz invariance is no longer manifest; we return to this in
\cref{sec:quantum:lorentz}. The field $X^+$ solves the wave equation, so
$X^+=X^+_L(\sigma^+)+X^+_R(\sigma^-)$, and the residual freedom of \cref{eq:quantum:residual} is
exactly a choice of what the functions $X_L^+$ and $X_R^+$ are. One uses it to remove all
oscillators from $X^+$:
\begin{equation}\label{eq:quantum:lcgauge}
 X^+=x^++\ap p^+\tau\qquad\text{(lightcone gauge)}
\end{equation}
\index{lightcone gauge}\source{Tong §2.2.1, ll.~2036--2066}. The zero modes $x^+$ and $p^+$
survive, because a reparametrization must respect the periodicity in $\sigma$. In this
gauge the worldsheet time $\tau$ is, up to a shift and a scale, the spacetime coordinate
$X^+$: every point of the string carries the same ``lightcone clock''. The gauge is good as long as $p^+\neq0$.

\subsection{Solving for \texorpdfstring{$X^-$}{X-}}

The benefit of the gauge choice in \cref{eq:quantum:lcgauge} is that the constraints become \emph{linear} in
$X^-$. Written with \cref{eq:quantum:lc}, the constraint $(\partial_+X)^2=0$ is
$2\,\partial_+X^-\partial_+X^+=\sum_i\partial_+X^i\partial_+X^i$, and
\cref{eq:quantum:lcgauge} gives $\partial_+X^+=\frac12\ap p^+$. Hence
\begin{equation}\label{eq:quantum:Xminus}
 \partial_+X^-=\frac{1}{\ap p^+}\sum_{i=1}^{D-2}\partial_+X^i\partial_+X^i,\qquad
 \partial_-X^-=\frac{1}{\ap p^+}\sum_{i=1}^{D-2}\partial_-X^i\partial_-X^i
\end{equation}
\source{Tong §2.2.1, ll.~2068--2113}. The field $X^-$ is determined by the transverse fields
up to one integration constant $x^-$. In modes, insert
$\partial_-X^-=\sqrt{\ap/2}\sum_n\alpha^-_n\ee^{-\ii n\sigma^-}$ and
$(\partial_-X^i)^2=\frac\ap2\sum_n\sum_m\alpha^i_{n-m}\alpha^i_m\ee^{-\ii n\sigma^-}$ and
compare coefficients:
\begin{equation}\label{eq:quantum:alphaminus}
 \alpha^-_n=\frac{1}{\sqrt{2\ap}\,p^+}\sum_{m\in\Z}\sum_{i=1}^{D-2}\alpha^i_{n-m}\alpha^i_m
\end{equation}
\source{Tong §2.2.1, ll.~2114--2141}. The case $n=0$, with $\alpha_0^-=\sqrt{\ap/2}\,p^-$ and
$\alpha_0^i\alpha_0^i=\frac\ap2p^ip^i$, is an equation for $p^-$:
\begin{equation}\label{eq:quantum:pminus}
 \frac{\ap p^-}{2}=\frac{1}{2p^+}\sum_{i=1}^{D-2}
   \Big(\frac\ap2\,p^ip^i+\sum_{n\neq0}\alpha^i_{-n}\alpha^i_n\Big),
\end{equation}
and the left-movers give the same equation with $\tilde\alpha$
\source{Tong §2.2.1, ll.~2142--2200}. Now compute the mass. Since
$p^2=-2p^+p^-+p^ip^i$,
\begin{equation}\label{eq:quantum:lcmass}
 M^2=2p^+p^--p^ip^i
    =\frac2\ap\sum_i\sum_{n\ne0}\alpha^i_{-n}\alpha^i_n
    =\frac4\ap\sum_{i=1}^{D-2}\sum_{n>0}\alpha^i_{-n}\alpha^i_n
    =\frac4\ap\sum_{i=1}^{D-2}\sum_{n>0}\tilde\alpha^i_{-n}\tilde\alpha^i_n ,
\end{equation}
where the third step uses that classical amplitudes commute. This is
\cref{eq:quantum:classmass} again, with the essential difference that only the $D-2$
transverse directions appear \source{Tong §2.2.1, ll.~2201--2238}. The independent variables
are $x^i,p^i,x^-,p^+$ and the transverse oscillators $\alpha^i_n,\tilde\alpha^i_n$. Note that
$p^-$ generates translations of $x^+$, that is, of $\tau$: \cref{eq:quantum:pminus} is the
lightcone Hamiltonian of the string.

%=====================================================================================
\section{Lightcone quantization and the normal-ordering constant}\label{sec:quantum:a}

\subsection{A positive Hilbert space}

Quantize the independent variables:
\begin{equation}\label{eq:quantum:lcccr}
 [x^i,p^j]=\ii\delta^{ij},\qquad [x^-,p^+]=-\ii,\qquad
 [\alpha^i_n,\alpha^j_m]=[\tilde\alpha^i_n,\tilde\alpha^j_m]=n\,\delta^{ij}\delta_{n+m,0}
\end{equation}
\source{Tong §2.2.2, ll.~2239--2256}. (That the naive commutators survive the elimination of
$X^\pm$ requires an argument, because $\alpha^-_n$ is a function of the $\alpha^i_n$. It
holds because $X^+$ in \cref{eq:quantum:lcgauge} has no oscillators, so the symplectic form
receives no contribution from $\alpha_n^-$ \source{Tong §2.2.2, ll.~2294--2313}.) The Fock
space is built on $\ket{0;p}$ with the transverse creation operators only. Since
$\delta^{ij}$ is positive definite, every state has positive norm. There are no ghosts, by
construction \source{Tong §2.2.2, ll.~2276--2290}.

The constraint of \cref{eq:quantum:pminus} must be imposed on states, and it has the ordering
ambiguity met before. With the \emph{level operators}\index{level (oscillator number)}
\begin{equation}\label{eq:quantum:N}
 N=\sum_{i=1}^{D-2}\sum_{n>0}\alpha^i_{-n}\alpha^i_n,\qquad
 \tilde N=\sum_{i=1}^{D-2}\sum_{n>0}\tilde\alpha^i_{-n}\tilde\alpha^i_n ,
\end{equation}
the quantum mass formula is
\begin{equation}\label{eq:quantum:massa}
 M^2=\frac4\ap(N-a)=\frac4\ap(\tilde N-a)
\end{equation}
\source{Tong §2.2.2, ll.~2318--2387}. In terms of \cref{eq:quantum:ladder},
$\alpha_{-n}\alpha_n=n\,a_n^\dagger a_n$, so $N$ is not the number of quanta but the sum of
their mode numbers: $\alpha^i_{-2}\ket{0;p}$ and $\alpha^i_{-1}\alpha^j_{-1}\ket{0;p}$ both have
$N=2$. Indeed, from $[\alpha^i_{-m}\alpha^i_m,\alpha^j_{-n}]=n\,\alpha^j_{-n}\delta_{mn}$ one
finds $[N,\alpha^j_{-n}]=n\,\alpha^j_{-n}$: each $\alpha_{-n}$ raises $N$ by $n$.

\subsection{Level matching as a symmetry}\label{sec:quantum:levelmatching}

The second equality in \cref{eq:quantum:massa}, $N=\tilde N$, has a simple meaning, which
we derive because it will be modified in \cref{ch:circle}. No point on a closed string is
marked, so a rigid shift $\sigma\to\sigma+\delta$ must act trivially on physical states.
In \cref{eq:quantum:modes} the shift sends $\ee^{-\ii n\sigma^-}\to\ee^{\ii n\delta}
\ee^{-\ii n\sigma^-}$ and $\ee^{-\ii n\sigma^+}\to\ee^{-\ii n\delta}\ee^{-\ii n\sigma^+}$, that
is,
\[
 \alpha_n\to\ee^{\ii n\delta}\alpha_n,\qquad \tilde\alpha_n\to\ee^{-\ii n\delta}\tilde\alpha_n .
\]
Since $[N,\alpha_n]=-n\alpha_n$, the operator $U_\delta=\ee^{-\ii\delta(N-\tilde N)}$
implements exactly this: $U_\delta\alpha_nU_\delta^{-1}=\ee^{\ii n\delta}\alpha_n$ and
$U_\delta\tilde\alpha_nU_\delta^{-1}=\ee^{-\ii n\delta}\tilde\alpha_n$. A state is invariant
under all $U_\delta$ if and only if $(N-\tilde N)\ket\psi=0$. Level matching is the
statement that the origin of $\sigma$ is unobservable. On a circle the zero modes also
transform, and the condition becomes $N-\tilde N=nw$ (\cref{ch:circle}).

\subsection{The constant \texorpdfstring{$a$}{a} as a Casimir energy}

Here is Tong's heuristic determination of $a$ \source{Tong §2.2.2, ll.~2391--2498}. Take the
classical expression $\frac12\sum_{n\neq0}\alpha^i_{-n}\alpha^i_n$ from
\cref{eq:quantum:pminus} at face value and bring it to normal-ordered form. The terms with
$n<0$ are in the wrong order; relabel $n\to-n$ and use
$\alpha^i_n\alpha^i_{-n}=\alpha^i_{-n}\alpha^i_n+n(D-2)$ (summed over $i$):
\begin{equation}\label{eq:quantum:ordering}
 \frac12\sum_{n\neq0}\alpha^i_{-n}\alpha^i_n
  =\frac12\sum_{n>0}\alpha^i_{-n}\alpha^i_n+\frac12\sum_{n>0}\alpha^i_{n}\alpha^i_{-n}
  =\sum_{n>0}\alpha^i_{-n}\alpha^i_n+\frac{D-2}{2}\sum_{n>0}n .
\end{equation}
The last term is the sum of the zero-point energies $\frac12\omega_n=\frac12n$ of $D-2$ towers
of oscillators. It diverges, as the vacuum energy does in any field theory. In a finite
box, however, the \emph{finite part} of such a sum is physical: it is the Casimir
energy\index{Casimir energy}. Regulate with a cutoff $\epsilon\ll1$ that suppresses high
frequencies,
\begin{equation}\label{eq:quantum:cutoff}
 \sum_{n=1}^\infty n\,\ee^{-\epsilon n}
 =-\frac{\partial}{\partial\epsilon}\sum_{n=1}^\infty\ee^{-\epsilon n}
 =-\frac{\partial}{\partial\epsilon}\frac{1}{\ee^{\epsilon}-1}
 =\frac{1}{\epsilon^2}-\frac1{12}+\frac{\epsilon^2}{240}+O(\epsilon^4).
\end{equation}
The expansion follows from $1/(\ee^\epsilon-1)=1/\epsilon-\frac12+\frac{\epsilon}{12}
-\frac{\epsilon^3}{720}+\dots$; we have confirmed it with a symbolic check in sympy. The
divergent term $1/\epsilon^2$ is independent of the state and has the form of a
cosmological constant on the worldsheet. It must be cancelled by a counterterm, because such
a term would violate the Weyl invariance of the Polyakov action. The finite remainder is
\begin{equation}\label{eq:quantum:zeta}
 \sum_{n=1}^\infty n\ \longrightarrow\ -\frac1{12}
 \qquad\Longrightarrow\qquad a=\frac{D-2}{24},\qquad
 M^2=\frac4\ap\Big(N-\frac{D-2}{24}\Big).
\end{equation}
The same number is obtained from \emph{zeta-function regularization}%
\index{zeta-function regularization}: the series $\zeta(s)=\sum_{n\ge1}n^{-s}$ converges for
$\mathrm{Re}\,s>1$, has a unique analytic continuation, and $\zeta(-1)=-\frac1{12}$
\source{Tong §2.2.2, ll.~2507--2530}.

\begin{example}[The cutoff sum with numbers]\label{ex:quantum:casimir}
Take $\epsilon=0.1$. Summing the series numerically gives
$\sum_nn\,\ee^{-0.1n}=99.916708\dots$, while $1/\epsilon^2-\frac1{12}=99.916667$ and the next
term $\epsilon^2/240=0.000042$ accounts for the difference. With $\epsilon=0.01$ the sum is
$9999.9166671$ against $10^4-\frac1{12}=9999.9166667$. The divergent part changes by a factor
$100$; the constant $-\frac1{12}=-0.08333\dots$ does not move. This stability is what makes
the finite part meaningful.
\end{example}

\begin{pitfallbox}[Common pitfall: ``$1+2+3+\dots=-\frac1{12}$'']
The divergent series does not equal $-\frac1{12}$ in any ordinary sense, and nothing in the
argument uses such an equality. What is claimed is: (i) the regulated sum has the structure
(divergent, state-independent) $+$ (finite) $+$ (vanishing with the cutoff); (ii) the
divergent part is removed by a local counterterm that the symmetries require anyway; (iii)
the finite part does not depend on how the cutoff is chosen. Point (iii) is not proved by
\cref{eq:quantum:cutoff}, which uses one particular cutoff. The derivation that does not
discard infinities uses conformal field theory: a theory of central charge $c$ on a cylinder
of circumference $2\pi$ has ground-state energy $-(c+\tilde c)/24$, and a free boson has
$c=\tilde c=1$, so each transverse direction contributes $-\frac1{24}$ to $L_0$ and to
$\tilde L_0$ \source{Tong §4.4.1, ll.~5018--5074}; see \cref{ch:cft}.
\end{pitfallbox}

%=====================================================================================
\section{The spectrum}\label{sec:quantum:spectrum}

We can now list the states of the closed bosonic string level by level, using
$M^2=\frac4\ap\big(N-\frac{D-2}{24}\big)$ with $N=\tilde N$.

\subsection{The ground state: a tachyon}

The state $\ket{0;p}$ has $N=\tilde N=0$ and
\begin{equation}\label{eq:quantum:tachyon}
 M^2=-\frac{1}{\ap}\,\frac{D-2}{6}
\end{equation}
\source{Tong §2.3.1, ll.~2535--2544}. A particle with $M^2<0$ is a
\emph{tachyon}\index{tachyon}. The right reading is field-theoretic. If $T(X)$ is the
spacetime field whose quantum this particle is, then $M^2=\partial^2V/\partial T^2$ at
$T=0$, so $M^2<0$ says that $T=0$ is a \emph{maximum} of the potential $V(T)$
(\cref{fig:quantum:tachyon}). The expansion is carried out around an unstable point, just as
the Higgs field at $H=0$ is tachyonic. Whether $V(T)$ for the closed bosonic string has a
stable minimum is not known: the cubic term produces one, the quartic term removes it again,
and $T$ mixes with other scalars \source{Tong §2.3.1, ll.~2546--2580}. The superstring has no
tachyon, and the lessons of the bosonic string carry over to it; but the instability of the
bosonic vacuum should be kept in mind throughout Part~I.

\begin{figure}[t]
\centering
\begin{tikzpicture}[scale=0.95]
 % left: tachyon potential
 \begin{scope}
 \draw[-{Stealth}] (-2.4,0)--(2.8,0) node[above]{$T$};
 \draw[-{Stealth}] (0,-1.5)--(0,1.6) node[left]{$V(T)$};
 \draw[thick,domain=-1.9:1.5,samples=60] plot (\x,{0.9-0.5*\x*\x});
 \draw[thick,dashed,domain=0:0.95,samples=20]
   plot ({1.5+\x},{-0.225-1.5*\x+1.8*\x*\x});
 \fill (0,0.9) circle (1.6pt);
 \node[above right] at (0,0.9) {\small $T=0$};
 \node at (1.95,-0.85) {\small ?};
 \end{scope}
 % right: spectrum
 \begin{scope}[xshift=5.2cm,yshift=-0.9cm]
 \draw[-{Stealth}] (-0.3,0.8)--(5.2,0.8) node[below]{$N=\tilde N$};
 \draw[-{Stealth}] (0,-0.3)--(0,2.9) node[left]{$\ap M^2$};
 \foreach \n/\y/\lab in {0/0/-4,1/0.8/0,2/1.6/4,3/2.4/8}{
   \draw (-0.07,\y)--(0.07,\y); \node[left] at (-0.05,\y) {\small $\lab$};
   \fill (1.2*\n+0.6,\y) circle (2pt);
   \draw (1.2*\n+0.6,0.73)--(1.2*\n+0.6,0.87);
 }
 \draw[dotted] (0.6,0)--(4.2,2.4);
 \node[below] at (0.6,0.0) {\small $1$};
 \node[below right] at (1.8,0.8) {\small $24^2$};
 \node[below right] at (3.0,1.6) {\small $324^2$};
 \node[below right] at (4.2,2.4) {\small $3200^2$};
 \end{scope}
\end{tikzpicture}
\caption{Left: a negative $M^2$ means that the expansion point $T=0$ is a maximum of the
tachyon potential; whether a minimum exists elsewhere (dashed) is not known. Right: the first four levels of the
closed bosonic string in $D=26$, with $\ap M^2=4(N-1)$ and the number of states at each
level (\cref{ex:quantum:counting}).}
\label{fig:quantum:tachyon}
\end{figure}

\subsection{The first excited level and \texorpdfstring{$D=26$}{D=26}}

Level matching forbids a single oscillator. The first excited states are
\begin{equation}\label{eq:quantum:level1}
 \alpha^i_{-1}\tilde\alpha^j_{-1}\ket{0;p},\qquad i,j=1,\dots,D-2,\qquad
 M^2=\frac4\ap\Big(1-\frac{D-2}{24}\Big),
\end{equation}
$(D-2)^2$ states in all \source{Tong §2.3.2, ll.~2595--2622}. They transform as a two-index
tensor of the rotation group $SO(D-2)$ of the transverse directions, which is all that
lightcone gauge leaves manifest. They must nevertheless form a representation of the full
Lorentz group $SO(1,D-1)$. Wigner's classification\index{little group} says how
\source{Tong §2.3.2, ll.~2631--2650}:
\begin{itemize}[leftmargin=*]
\item a \emph{massive} particle can be brought to rest, $p^\mu=(M,0,\dots,0)$; the rotations
that preserve this momentum form $SO(D-1)$, and the internal states must fill
representations of that group, $SO(D-1)$;
\item a \emph{massless} particle cannot be brought to rest; one takes $p^\mu=(k,0,\dots,0,k)$,
which is preserved by $SO(D-2)$, and the internal states need only fill representations of
$SO(D-2)$. This is why a photon in four dimensions has two polarizations while a massive
vector has three.
\end{itemize}
The $(D-2)^2$ states of \cref{eq:quantum:level1} form a tensor of $SO(D-2)$. They cannot be
completed into a representation of $SO(D-1)$: a two-index tensor of $SO(D-1)$ would need
$(D-1)^2$ states, and nothing else is available at this level. Lorentz invariance therefore
requires these states to be massless, and \cref{eq:quantum:level1} then gives
\begin{equation}\label{eq:quantum:D26}
 1-\frac{D-2}{24}=0\qquad\Longleftrightarrow\qquad D=26,\qquad a=1
\end{equation}
\index{critical dimension}\TL\ \source{Tong §2.3.2, ll.~2643--2653}. This agrees with
\cref{ex:quantum:level1}: $a=1$ is the value at which the first level is a massless gauge
particle. With $D=26$ the mass formula takes the form used in the rest of the book,
\begin{equation}\label{eq:quantum:massfinal}
 M^2=\frac4\ap(N-1)=\frac{2}{\ap}\big(N+\tilde N-2\big),\qquad N=\tilde N .
\end{equation}

\subsection{Graviton, \texorpdfstring{$B$}{B}-field and dilaton}

For $D=26$ the massless states are a tensor $\zeta_{ij}$ of $SO(24)$ with $24^2=576$
components. A two-index tensor decomposes into irreducible parts,
\begin{equation}\label{eq:quantum:decomp}
 \underbrace{\mathbf{24}\otimes\mathbf{24}}_{576}
 =\underbrace{\text{traceless symmetric}}_{\frac12\cdot24\cdot25-1\,=\,299}
 \ \oplus\ \underbrace{\text{antisymmetric}}_{\frac12\cdot24\cdot23\,=\,276}
 \ \oplus\ \underbrace{\text{trace}}_{1},
\end{equation}
and to each part corresponds a massless field in spacetime
\source{Tong §2.3.2, ll.~2655--2684}:
\begin{center}
\begin{tabular}{@{}llll@{}}\toprule
$SO(24)$ representation & states & field & name\\\midrule
traceless symmetric & 299 & $G_{\mu\nu}(X)$ & metric; its quantum is the graviton\index{graviton}\\
antisymmetric & 276 & $B_{\mu\nu}(X)$ & Kalb--Ramond field, or 2-form\index{Kalb--Ramond field}\\
singlet & 1 & $\Phi(X)$ & dilaton\index{dilaton}\\\bottomrule
\end{tabular}
\end{center}
These three fields are common to all string theories. The fields carry spacetime indices
$\mu,\nu=0,\dots,25$ and not just $i,j$, because $G_{\mu\nu}$ and $B_{\mu\nu}$ possess gauge
symmetries that remove the extra components; lightcone gauge shows only the physical
polarizations, and \cref{ex:quantum:level1} showed where the gauge symmetry comes from
\source{Tong §2.3.2, ll.~2760--2770}.

\begin{physicsbox}[Why a massless spin-two particle means gravity]
A massless spin-two field $h_{\mu\nu}$ quantized covariantly has negative-norm states,
created by operators with one timelike index, for the same reason as
\cref{eq:quantum:negnorm}. They decouple only if the theory has the gauge symmetry
$h_{\mu\nu}\to h_{\mu\nu}+\partial_\mu\xi_\nu+\partial_\nu\xi_\mu$. An argument due to Feynman
and Weinberg shows that this symmetry survives interactions only if the full theory is
invariant under diffeomorphisms, that is, if it is general relativity, possibly with
higher-derivative corrections \source{Tong §2.3.2, ll.~2686--2757}. The string was not
designed to contain gravity. The graviton appears because the constraint algebra forces
$a=1$. How Einstein's equations themselves emerge is the subject of \cref{ch:backgrounds};
the roles of $B_{\mu\nu}$ and $\Phi$ in T-duality begin in \cref{ch:buscher}.
\end{physicsbox}

\subsection{Massive levels}

At higher levels all states are massive and must fill representations of $SO(D-1)=SO(25)$,
although only $SO(24)$ is manifest. At level $N=2$ the right-moving states are
$\alpha^i_{-1}\alpha^j_{-1}\ket0$ (symmetric in $i,j$) and $\alpha^i_{-2}\ket0$. Their number is
\begin{equation}\label{eq:quantum:level2}
 \tfrac12(D-2)(D-1)+(D-2)=\tfrac12D(D-1)-1 ,
\end{equation}
which is exactly the dimension of the traceless symmetric tensor of $SO(D-1)$
\source{Tong §2.3.3, ll.~2771--2806}. (A symmetric tensor of $SO(n)$ has $\frac12n(n+1)$
components, and removing the trace gives $\frac12n(n+1)-1$; put $n=D-1$.) The vector
$\alpha^i_{-2}\ket0$ of $SO(24)$ supplies precisely the components that the symmetric tensor of
$SO(24)$ lacks. One can show that all higher levels also assemble into $SO(D-1)$
representations, so the only condition from this argument is the one at level 1
\source{Tong §2.3.3, ll.~2807--2812}.

\begin{example}[Counting states with a generating function]\label{ex:quantum:counting}
Each right-moving oscillator $\alpha^i_{-n}$ may be excited $k=0,1,2,\dots$ times and then
contributes $kn$ to $N$. The number $d_N$ of right-moving states at level $N$ is therefore
the coefficient of $q^N$ in
\[
 \prod_{n=1}^\infty\Big(\frac{1}{1-q^n}\Big)^{24}=1+24\,q+324\,q^2+3200\,q^3+25650\,q^4+\dots
\]
(expansion obtained by a symbolic check in sympy). Check the first entries by hand:
$d_1=24$; $d_2=\frac12\cdot24\cdot25+24=324=\frac12\cdot25\cdot26-1$, in agreement with
\cref{eq:quantum:level2}; and $d_3=24+24^2+\binom{26}{3}=24+576+2600=3200$ from
$\alpha_{-3}$, $\alpha_{-2}\alpha_{-1}$ and $\alpha_{-1}\alpha_{-1}\alpha_{-1}$. For $SO(25)$
the traceless symmetric three-index tensor has $\binom{27}{3}-25=2900$ components and the
antisymmetric two-index tensor has $\binom{25}{2}=300$; the sum is $3200$, as required for a
massive level. The closed string at level $N=\tilde N$ has $d_N^2$ states:
\begin{center}
\begin{tabular}{@{}ccccl@{}}\toprule
$N=\tilde N$ & $\ap M^2$ & $d_N$ & closed-string states & content\\\midrule
0 & $-4$ & 1 & 1 & tachyon\\
1 & $0$ & 24 & 576 & $G_{\mu\nu}$, $B_{\mu\nu}$, $\Phi$\\
2 & $4$ & 324 & 104\,976 & massive, from $\mathbf{324}\otimes\mathbf{324}$ of $SO(25)$\\
3 & $8$ & 3200 & 10\,240\,000 & massive\\\bottomrule
\end{tabular}
\end{center}
The spacing of the levels is $\Delta M^2=4/\ap$. If $\sqrt\ap$ is of the order of the Planck
length, the first massive level has $M=2/\sqrt{\ap}$, near the Planck mass, far beyond any
accelerator \source{Tong §2.3.3, ll.~2813--2819}. At low energy only the massless level
matters; this is why \cref{ch:backgrounds} can describe the string by a field theory of
$G$, $B$ and $\Phi$.
\end{example}

The same product appears, with $q=\ee^{2\pi\ii\tau}$, in the one-loop partition function of the
string, and its analogue for K3 in the elliptic genus of \cref{ch:ellgenus}.

%=====================================================================================
\section{Lorentz invariance revisited}\label{sec:quantum:lorentz}

The argument for $D=26$ in the previous section was quick: it fixed $a$ by a regularization
and $D$ by one representation-theoretic observation. A more systematic test within lightcone
quantization is to construct the Lorentz generators and verify their algebra
\source{Tong §2.4, ll.~2820--2835}.

Before gauge fixing, the action is invariant under the global Poincar\'e transformations
$X^\mu\to\Lambda^\mu{}_\nu X^\nu+c^\mu$. The Noether currents are
\begin{equation}\label{eq:quantum:currents}
 P^\alpha_\mu=T\,\partial^\alpha X_\mu,\qquad
 J^\alpha_{\mu\nu}=P^\alpha_\mu X_\nu-P^\alpha_\nu X_\mu ,
\end{equation}
conserved by the equations of motion \source{Tong §2.4, ll.~2836--2860}. The charges
$M^{\mu\nu}=\int_0^{2\pi}\dd\sigma\,J^{\tau\,\mu\nu}$ are, in modes,
\begin{equation}\label{eq:quantum:Mmunu}
 M^{\mu\nu}=\underbrace{(p^\mu x^\nu-p^\nu x^\mu)}_{l^{\mu\nu}}
  \ \underbrace{-\,\ii\sum_{n=1}^\infty\frac1n\big(\alpha^\nu_{-n}\alpha^\mu_n
      -\alpha^\mu_{-n}\alpha^\nu_n\big)}_{S^{\mu\nu}}
  \ +\ \tilde S^{\mu\nu},
\end{equation}
where $l^{\mu\nu}$ is the orbital angular momentum of the centre of mass, $S^{\mu\nu}$ the
spin carried by right-moving oscillators and $\tilde S^{\mu\nu}$ the same expression with
$\tilde\alpha$ \source{Tong §2.4, ll.~2861--2880}. In covariant quantization these operators
obey the Lorentz algebra\index{Lorentz algebra}
\begin{equation}\label{eq:quantum:lorentzalg}
 [M^{\rho\sigma},M^{\tau\nu}]=\eta^{\sigma\tau}M^{\rho\nu}-\eta^{\rho\tau}M^{\sigma\nu}
   +\eta^{\rho\nu}M^{\sigma\tau}-\eta^{\sigma\nu}M^{\rho\tau}
\end{equation}
without difficulty. In lightcone gauge the generators $M^{i-}$ are dangerous. They contain
$\alpha^-_n$ and $p^-$, which by \cref{eq:quantum:alphaminus,eq:quantum:pminus} are
\emph{quadratic} in the transverse oscillators and carry the ordering constant $a$. So
$M^{i-}$ is cubic in oscillators, and its commutators can produce anomalous terms. The
algebra of \cref{eq:quantum:lorentzalg} demands
\begin{equation}\label{eq:quantum:smoking}
 [M^{i-},M^{j-}]=0 ,
\end{equation}
because with $\rho=i$, $\tau=j$, $\sigma=\nu=-$ every term on the right of
\cref{eq:quantum:lorentzalg} vanishes: $\eta^{--}=\eta^{i-}=\eta^{-j}=0$ in the metric
of \cref{eq:quantum:lc}, and the remaining term is $-\delta^{ij}M^{--}$, which is zero by
antisymmetry. Keeping $a$ and $D$ arbitrary, a long calculation gives
instead
\begin{equation}\label{eq:quantum:anomaly}
 [M^{i-},M^{j-}]=\frac{2}{(p^+)^2}\sum_{n>0}
  \Big(\Big[\frac{D-2}{24}-1\Big]n+\frac1n\Big[a-\frac{D-2}{24}\Big]\Big)
  \big(\alpha^i_{-n}\alpha^j_n-\alpha^j_{-n}\alpha^i_n\big)+(\alpha\leftrightarrow\tilde\alpha)
\end{equation}
\TL\ \source{Tong §2.4, ll.~2905--2945}. The operators
$\alpha^i_{-n}\alpha^j_n-\alpha^j_{-n}\alpha^i_n$ for different $n$ are independent, so each
coefficient must vanish. The coefficient is a combination of $n$ and $1/n$, and it vanishes
for all $n\geq1$ only if both brackets vanish:
\begin{equation}\label{eq:quantum:D26a1}
 \frac{D-2}{24}=1\quad\text{and}\quad a=\frac{D-2}{24}
 \qquad\Longleftrightarrow\qquad D=26,\quad a=1
\end{equation}
\source{Tong §2.4, ll.~2953--2956}. Two remarks. First, the second condition is
\cref{eq:quantum:zeta}: the Lorentz algebra confirms the Casimir value of $a$ without any
regularization. Second, this is an \emph{anomaly}\index{anomaly}: a symmetry of the classical
theory that fails in the quantum theory unless a condition is met. We have not carried out
the calculation leading to \cref{eq:quantum:anomaly}; neither does Tong, and neither does
the Lean library (\cref{sec:quantum:lean}).

\begin{historybox}
Lightcone quantization of the string was first carried out by Goddard, Goldstone, Rebbi and
Thorn, ``Quantum Dynamics of a Massless Relativistic String'', Nucl.\ Phys.\ B56 (1973); the
commutator in \cref{eq:quantum:anomaly} is the central computation of that paper
\source{Tong §2.4, footnote, ll.~2949--2951}.
\end{historybox}

%=====================================================================================
\section{The critical dimension as a central charge; the superstring}
\label{sec:quantum:centralcharge}

The number 26 will be derived a third time in \cref{ch:ghosts}, and that derivation is the
one formalized in Lean, so we summarize it here. Gauge-fixing the Polyakov path integral
introduces anticommuting Faddeev--Popov fields $b,c$. They form a conformal field theory, and
any conformal field theory has a \emph{central charge}\index{central charge} $c$, a number
that measures its response to a rescaling of the worldsheet metric (\cref{ch:cft}). The
results needed are \TL\ \source{Tong §5.2--5.3, ll.~6832--6850}:
\begin{itemize}[leftmargin=*]
\item a free scalar field has $c=1$, so $D$ fields $X^\mu$ have $c=D$;
\item the $bc$ ghost system has $c=-26$;
\item the Weyl symmetry is anomalous unless the total central charge vanishes.
\end{itemize}
Since the Weyl symmetry is a gauge symmetry, it must not be anomalous. Hence $D-26=0$. In
this form the statement is more general than ``spacetime has 26 dimensions'': the matter
sector may be \emph{any} conformal field theory with $c=26$, for instance four free scalars
for Minkowski space together with an ``internal'' theory of $c=22$. The critical dimension
is really a critical central charge \source{Tong §5.3, ll.~6851--6868}.

The lightcone result and the central-charge result are connected by the Casimir energy. A
conformal field theory on the cylinder has ground-state energy $-(c+\tilde c)/24$
\source{Tong §4.4.1, ll.~5060--5074}, which gives $a=c_\perp/24$ for a transverse theory of
central charge $c_\perp$; with $c_\perp=D-2$ this is \cref{eq:quantum:zeta}. The identity
\[
 a=\frac{D-2}{24}=1+\frac{D-26}{24}
\]
then shows that the ordering constant takes the ghost-free value $a=1$ of
\cref{ex:quantum:level1} exactly when the total central charge $D-26$ vanishes. In both
languages two of the $D$ directions do not contribute physical oscillators: in lightcone
gauge they are eliminated by hand, in the covariant path integral the $bc$ system cancels
them.

\paragraph{The superstring.} Adding fermions $\psi^\mu$ to the worldsheet, related to the
$X^\mu$ by worldsheet supersymmetry, gives the superstring. Supersymmetry is a further gauge
symmetry and requires further Faddeev--Popov fields, the commuting $\beta\gamma$ system, with
$c=+11$. The matter sector must then have $c=26-11=15$. A free fermion has $c=\frac12$, and
supersymmetry pairs each boson with a fermion, so $D(1+\frac12)=15$ and
\begin{equation}\label{eq:quantum:D10}
 D=10
\end{equation}
\index{superstring}\TL\ \source{Tong §5.3.1, ll.~6884--6891}. The quantized superstring has no
tachyon, and its massless bosons include $G_{\mu\nu}$, $B_{\mu\nu}$ and $\Phi$ again
\source{Tong §2.5, ll.~2957--2975}. There are several superstring theories, distinguished by
discrete choices: type~II (worldsheet fermions moving in both directions, 32 supercharges in
$D=10$), subdivided into IIA and IIB by their additional Ramond--Ramond fields; heterotic
(fermions in one direction only, 16 supercharges) with gauge group $SO(32)$ or $E_8\times
E_8$; and type~I with open strings \source{Tong §2.5, ll.~2976--3008}. For this book the
important consequences are these. In $D=10$ one can compactify six dimensions, for instance
on $\KT$ (\cref{ch:k3t2}). By the same counting, a supersymmetric sigma model on a space of
four real dimensions, such as a K3 surface, has four bosons and four fermions and hence
$c=4\cdot1+4\cdot\frac12=6$, the value used in \cref{ch:ellgenus}. And the appearance of the lattice $E_8$ in a string
theory is one motivation for \cref{ch:lattices}.

%=====================================================================================
\section{What the machine has checked}\label{sec:quantum:lean}

The Lean file for this chapter is
\begin{center}
\path{StringTheoryFormalization/UseCases/CriticalDimension.lean}
\end{center}
and all its declarations live in the namespace
\lean{StringTheory.UseCases.CriticalDimension}. The file belongs to the older part
of the library, in which a physical quantity is modelled by a rational number and the
theorems are arithmetic. It is a good first Lean file to read, because every proof is one
line. It is also a good first lesson in reading a formal statement for what it literally
says.

\begin{leanbox}[In Lean: the two ghost central-charge formulas]
\begin{leancode}
def fermionicGhostCharge (lam : ℚ) : ℚ := 1 - 3 * (2 * lam - 1) ^ 2

def bosonicGhostCharge (lam : ℚ) : ℚ := 3 * (2 * lam - 1) ^ 2 - 1

theorem bosonic_eq_neg_fermionic (lam : ℚ) :
    bosonicGhostCharge lam = - fermionicGhostCharge lam := by ...
\end{leancode}
\textbf{Reading.} These are \emph{definitions} of two functions $\Q\to\Q$:
$c_{bc}(\lambda)=1-3(2\lambda-1)^2$ and $c_{\beta\gamma}(\lambda)=3(2\lambda-1)^2-1$. In the
literature they are the central charges of an anticommuting, respectively commuting,
first-order system whose fields have conformal weights $\lambda$ and $1-\lambda$ (the file's
header cites Polchinski and Green--Schwarz--Witten for them; these are not in this book's
source library; Tong computes only the case $\lambda=2$, where $b$ has weight $2$ and $c$
weight $-1$ \source{Tong §5.2, ll.~6700--6839}). Lean has not derived these
formulas from an operator product expansion; they are input. The theorem states that the
second polynomial is the negative of the first. \textbf{Proof idea.} Unfold the definitions;
the tactic \lean{ring} verifies the polynomial identity.
\end{leanbox}

\begin{leanbox}[In Lean: evaluation and the two cancellations]
\begin{leancode}
theorem bc_ghost_charge_eq : fermionicGhostCharge 2 = -26 := by ...

theorem betagamma_ghost_charge_eq : bosonicGhostCharge (3 / 2) = 11 := by ...

theorem bosonic_string_critical_dimension :
    (26 : ℚ) * 1 + fermionicGhostCharge 2 = 0 := by ...

theorem superstring_critical_dimension :
    (10 : ℚ) * 1 + 10 * (1 / 2) + fermionicGhostCharge 2
      + bosonicGhostCharge (3 / 2) = 0 := by ...
\end{leancode}
\textbf{Reading.} The first two theorems evaluate the polynomials: $1-3\cdot3^2=-26$ and
$3\cdot2^2-1=11$. The third says $26\cdot1+(-26)=0$ in $\Q$; the fourth says
$10\cdot1+10\cdot\frac12-26+11=0$ in $\Q$. The literals \texttt{1} and \texttt{1/2} stand for
the central charges of a free boson and a free fermion, but in the statement they are just
rational numbers. \textbf{Proof idea.} \lean{norm\_num} evaluates the arithmetic after the
definitions are unfolded; the last two theorems first rewrite with the first two.
\end{leanbox}

\begin{leanbox}[In Lean: uniqueness of the solution]
\begin{leancode}
theorem bosonic_dimension_unique (D : ℚ)
    (h : D * 1 + fermionicGhostCharge 2 = 0) :
    D = 26 := by ...
\end{leancode}
\textbf{Reading.} If a rational number $D$ satisfies $D\cdot1+c_{bc}(2)=0$, then $D=26$. This
is the solution of one linear equation. Note that $D$ is a rational number, not a natural
number, so the theorem is slightly stronger than ``no other integer dimension works''.
\textbf{Proof idea.} Rewrite the hypothesis with \lean{bc\_ghost\_charge\_eq} to obtain
$D-26=0$, then \lean{linarith}.
\end{leanbox}

\paragraph{What this does and does not establish.} The chain of reasoning behind $D=26$ in
the central-charge form has four links: (1) gauge fixing produces the $bc$ system with
$\lambda=2$; (2) a first-order system of weight $\lambda$ has $c=1-3(2\lambda-1)^2$, and a
free boson has $c=1$; (3) consistency requires $c_{\rm total}=0$; (4) arithmetic. The Lean
file checks link (4) and records link (2) as a definition. Links (1)--(3) are physics and
conformal field theory established in the literature. In particular the file does
\emph{not}:
\begin{itemize}[leftmargin=*]
\item define the Virasoro algebra, an operator product expansion or a central charge as a
mathematical object (so it cannot \emph{derive} $-26$);
\item contain the oscillators $\alpha^\mu_n$, the generators $M^{\mu\nu}$, or the commutator
of \cref{eq:quantum:anomaly}; it does not derive the closure of the Lorentz algebra, which is
the content of \cref{sec:quantum:lorentz};
\item determine the normal-ordering constant $a$, prove a no-ghost theorem, or construct
any conformal field theory with $c=26$.
\end{itemize}
In the language of this book, the file proves an \emph{arithmetic shadow} of the
critical-dimension statement. A faithful formalization would need, at least, a Lie-algebra
definition of the Virasoro algebra with its central extension, the Fock-space representation
built from \cref{eq:quantum:osc}, and a proof that the normal-ordered $L_n$ satisfy the
algebra with $c=D$; Exercise~\ref{exo:quantum:virasoro} indicates the first step. All
theorems of the file depend only on the axioms \lean{propext}, \lean{Classical.choice} and
\lean{Quot.sound}; this was checked with \lean{\#print axioms} for each of the six theorems.
\Cref{tab:quantum:tiers} collects the status of the claims of this chapter.

\begin{table}[t]
\centering\small
\caption{Epistemic status of the main claims of this chapter.}\label{tab:quantum:tiers}
\begin{tabular}{@{}>{\raggedright\arraybackslash}p{0.43\textwidth}cp{0.45\textwidth}@{}}\toprule
Claim & Tier & Basis\\\midrule
$1-3(2\lambda-1)^2=-26$ at $\lambda=2$ & \TA & \lean{bc\_ghost\_charge\_eq}\\
$3(2\lambda-1)^2-1=11$ at $\lambda=\frac32$ & \TA & \lean{betagamma\_ghost\_charge\_eq}\\
$26\cdot1-26=0$ in $\Q$ & \TA & \lean{bosonic\_string\_critical\_dimension}\\
$10\cdot1+10\cdot\frac12-26+11=0$ in $\Q$ & \TA & \lean{superstring\_critical\_dimension}\\
$D\in\Q$ and $D-26=0$ imply $D=26$ & \TA & \lean{bosonic\_dimension\_unique}\\
$c_{\beta\gamma}(\lambda)=-c_{bc}(\lambda)$ as polynomials & \TA & \lean{bosonic\_eq\_neg\_fermionic}\\
The $bc$ ghosts have $c=-26$, the $\beta\gamma$ ghosts $c=11$ & \TL & Tong §5.2--5.3.1\\
Weyl invariance needs $c_{\rm total}=0$, so $D=26$ (resp.\ $10$) & \TL & Tong §5.3, §5.3.1\\
Lightcone Lorentz algebra closes iff $D=26$, $a=1$ & \TL & Tong §2.4\\
$\sum_n n\,\ee^{-\epsilon n}=\epsilon^{-2}-\frac1{12}+O(\epsilon^2)$;
  $d_N=1,24,324,3200$ & --- & symbolic check (sympy), not Tier A\\
\bottomrule
\end{tabular}
\end{table}
The identification of the rational numbers in these theorems with central charges of
worldsheet theories is part of the physical modelling and carries no tier~A status.

%=====================================================================================
\begin{summarybox}
\begin{itemize}[leftmargin=*]
\item Canonical quantization turns the modes of $X^\mu$ into oscillators,
$[\alpha^\mu_n,\alpha^\nu_m]=n\eta^{\mu\nu}\delta_{n+m,0}$. The timelike oscillators create
negative-norm states, which the Virasoro constraints $L_n\ket{\rm phys}=0$ $(n>0)$ and
$(L_0-a)\ket{\rm phys}=0$ must remove.
\item In lightcone gauge $X^+=x^++\ap p^+\tau$ the constraints determine $X^-$; only $D-2$
transverse oscillators remain and the Hilbert space is positive. The price is that Lorentz
invariance must be checked.
\item The mass formula is $M^2=\frac4\ap(N-a)$ with $N=\tilde N$ (level matching $=$ invariance
under shifts of $\sigma$). The ordering constant is a Casimir energy, $a=(D-2)/24$.
\item Level 0 is a tachyon, $M^2=-4/\ap$ for $D=26$. Level 1 has $(D-2)^2$ states, which can
represent the Lorentz group only if massless: $D=26$, $a=1$. They are the graviton
$G_{\mu\nu}$, the Kalb--Ramond field $B_{\mu\nu}$ and the dilaton $\Phi$ ($299+276+1=576$).
\item The same values follow from closure of the lightcone Lorentz algebra,
$[M^{i-},M^{j-}]=0$, and from cancellation of the central charge, $D-26=0$. For the
superstring, $\frac32D-26+11=0$ gives $D=10$.
\item Lean checks the central-charge arithmetic \TA. The origin of the numbers $-26$, $11$,
$1$, $\frac12$ and the requirement $c_{\rm total}=0$ are \TL.
\end{itemize}
\end{summarybox}

\section*{Exercises}
\addcontentsline{toc}{section}{Exercises}

\begin{exercise}[$\star$]
Using \cref{eq:quantum:osc}, show that $[N,\alpha^j_{-n}]=n\,\alpha^j_{-n}$ and
$[N,\alpha^j_n]=-n\,\alpha^j_n$ for $n>0$. Deduce that
$\alpha^{i}_{-2}\alpha^{j}_{-1}\alpha^{k}_{-1}\ket{0;p}$ has $N=4$, and find the
state(s) with which level matching allows it to be combined if the left-moving part contains
only oscillators $\tilde\alpha_{-1}$ and $\tilde\alpha_{-3}$.
\end{exercise}

\begin{exercise}[$\star$]
Compute the norm of $\alpha^0_{-1}\alpha^0_{-1}\ket{0;p}$ and of
$\alpha^0_{-1}\alpha^1_{-1}\ket{0;p}$ (strip off the momentum delta function). Conclude that
the sign of the norm is $(-1)^{k}$, with $k$ the number of timelike oscillators.
\end{exercise}

\begin{exercise}[$\star$]
Verify the decomposition in \cref{eq:quantum:decomp} for general $n=D-2$: show that
$\big[\frac12n(n+1)-1\big]+\frac12n(n-1)+1=n^2$. In $D=4$ (not a consistent bosonic string,
but a useful check) how many polarizations would the ``graviton'', ``$B$-field'' and
``dilaton'' have? Compare with the two helicities of the four-dimensional graviton.
\end{exercise}

\begin{exercise}[$\star\star$]
Repeat the regularization of \cref{eq:quantum:cutoff} with a Gaussian cutoff, that is, for
the sum $\sum_nn\,\ee^{-\epsilon^2n^2}$. Use the Euler--Maclaurin formula
\[
 \sum_{n\ge0}f(n)=\int_0^\infty f(x)\,\dd x+\tfrac12f(0)-\tfrac1{12}f'(0)+\tfrac1{720}f'''(0)-\dots
\]
Show that the
divergent part changes but the finite part is again $-\frac1{12}$. Why does every smooth
cutoff function $f(n)=n\,g(\epsilon n)$ with $g(0)=1$ give the same finite part?
\end{exercise}

\begin{exercise}[$\star\star$]
Open strings with free ends in all directions (\cref{ch:dbranes}) have a single set of
oscillators and $M^2=\frac1\ap(N-a)$ \source{Tong §3.1, ll.~3270--3290}, with the same
$a=(D-2)/24$. Repeat the argument of
\cref{sec:quantum:spectrum}: how many states are there at $N=1$, which little group must they
represent, and what follows for $D$? Identify the massless particle.
\end{exercise}

\begin{exercise}[$\star\star$]
(a) Show that the number of right-moving states at level 4 in $D=26$ is $25650$ by listing
the partitions of 4 and counting the states for each. (b) For $SO(25)$ the traceless
symmetric tensors of rank $r$ have dimension $\binom{24+r}{r}-\binom{22+r}{r-2}$, and the
traceless three-index tensor of mixed symmetry has dimension $\frac13n(n^2-1)-n$ with
$n=25$. Check that the traceless symmetric tensors of rank 4 and rank 2, the mixed tensor
and a singlet have total dimension $25650$. (A dimension count does not prove that the states
transform in these representations; what else would be needed?)
\end{exercise}

\begin{exercise}[$\star\star$]\label{exo:quantum:superunique}
The library proves uniqueness only for the bosonic string. State and prove in Lean the
superstring analogue: for a rational number $D$, the hypothesis
\begin{leancode}
D * 1 + D * (1 / 2) + fermionicGhostCharge 2 + bosonicGhostCharge (3 / 2) = 0
\end{leancode}
implies $D=10$. Hint: rewrite the hypothesis with the two evaluation theorems of
\cref{sec:quantum:lean}, then use \lean{linarith}. This statement is an exercise; it is not
a declaration of the library.
\end{exercise}

\begin{exercise}[$\star\star$]
Evaluate \lean{fermionicGhostCharge} at $\lambda=1$ and $\lambda=\frac12$ by hand, and prove
both values in Lean with \lean{norm\_num}. Then prove that \lean{fermionicGhostCharge} is
symmetric under $\lambda\to1-\lambda$, i.e.\
\lean{fermionicGhostCharge (1 - lam) = fermionicGhostCharge lam}. What is the physical
meaning of this symmetry for a system with fields of weights $\lambda$ and $1-\lambda$?
\end{exercise}

\begin{exercise}[$\star\star\star$]\label{exo:quantum:virasoro}
(a) Using only \cref{eq:quantum:osc}, show that $[L_m,\alpha^\mu_n]=-n\,\alpha^\mu_{m+n}$.
(b) Deduce that $[L_m,L_n]=(m-n)L_{m+n}$ for $m+n\neq0$. (c) For $m+n=0$ the normal ordering
of $L_0$ produces an extra constant $A(m)$. Compute $A(1)$ and $A(2)$ by evaluating
$\bra{0;0}[L_m,L_{-m}]\ket{0;0}$ directly, and check that they agree with
$A(m)=\frac{D}{12}m(m^2-1)$. (d) Write a Lean \texttt{structure} whose fields are a family
$L:\Z\to V$ in a Lie algebra $V$ over $\Q$, a central element, and the commutation relation
with central term $\frac{c}{12}m(m^2-1)\delta_{m+n,0}$. This is the first step towards a
formalization in which ``$c=D$'' could be a theorem and not a definition.
\end{exercise}

\section*{Further reading}
\addcontentsline{toc}{section}{Further reading}
\begin{itemize}[leftmargin=*]
\item Covariant quantization, ghosts and the constraints:
\source{Tong §2.1, ll.~1708--1931}.
\item Lightcone gauge, quantization, the Casimir energy and zeta-function regularization:
\source{Tong §2.2, ll.~1932--2530}.
\item The spectrum, Wigner's little-group argument, and why massless spin two means general
relativity: \source{Tong §2.3, ll.~2533--2819}.
\item The Lorentz algebra in lightcone gauge and the original reference:
\source{Tong §2.4, ll.~2820--2956}. The superstring in outline:
\source{Tong §2.5, ll.~2957--3008}.
\item The Casimir energy from the Schwarzian derivative, without discarded infinities:
\source{Tong §4.4.1, ll.~5018--5074}. The ghost central charge and the critical central
charge: \source{Tong §5.2--5.3.1, ll.~6832--6891}; both are developed in
\cref{ch:cft,ch:ghosts}.
\item The Lean source, including its own account of what is and is not proved:\\
\path{StringTheoryFormalization/UseCases/CriticalDimension.lean}.
\end{itemize}
